% ============================================================
% 分片：s03c-applications  对应原文 PDF 第 29--34 页（印刷页码 29--34）
% 负责章节：3.3 Applications
% 译者：t06-s33   审校：t03-s2b（逐式核对完成）   审校状态：二审
% ============================================================
% 边界说明：
% [边界: 第29页顶部两段属 §3.2 收尾，置于本节标题前]
%   1) 第 29 页顶部两段为 §3.2 "The Singularity" 的收尾（第 28 页末段已完整，
%      无跨句分割），而 s03b 的分工止于第 28 页；按 Lead 裁定，这两段并入本片，
%      置于 \subsection{应用} 之前。这两段中出现的公式编号仍属 3.2（如 (3.16)），
%      照原文编号，不改号。
%   2) 第 34 页末的跨页半句（"The A model has a greatly enriched variant ...
%      vertex operators are"）按 Lead 裁定归 s03d，本片不收录。
% ============================================================

% ==== 原文 PDF 第 29 页（印刷页 29）====
\label{op:29}

人们相信，本文所研究的理论在红外都流向共形场论。若果真如此，那么从 Landau--Ginzburg 到 Calabi--Yau 的一种特别有趣的过渡方式，将是经由共形不变理论。是否存在一族连续的共形场论，从 Landau--Ginzburg 插值到 Calabi--Yau？这是 Landau--Ginzburg 与 Calabi--Yau 作为共形场论能够真正“等价”的唯一途径。在拓扑改变的情形（我们将在 §5.5 看到它与这里类似），镜像对称 [7] 直接给出了答案为“是”的迹象。在当前语境下，我们可以（不严格地）尝试如下论证。首先，一族共形场论若在实线上出现任何奇点或不连续，则它在圆周上必然已经出现，因为从共形角度看，实轴不过是挖去一点的圆周。其次，在上文中，对线性 $\sigma$ 模型我们已经找出了有限体积内出现的全部奇点：唯一的这类奇点位于 $r=\theta=0$。避免奇点的要点在于大 $\sigma$ 处最小能量密度 $U$ 的非零性，$U$ 由 (3.16) 给出。这是一个精确公式，不依赖于小 $e$。为趋近红外极限，应当令 $e\to\infty$；而看起来 $U$ 在该极限下只会增大，因此不会产生新的奇点。这一论证使作者倾向于相信：对有限体积的线性 $\sigma$ 模型所找到的奇点，就是共形场论中会出现的全部奇点。该论证潜在的麻烦在于，$U$ 的定义涉及大 $\sigma$ 极限，而这一极限未必与 $e\to\infty$ 极限交换。

% [Tutor Note: 这里的论证只针对“族在实轴上的奇点必在圆周上出现”这一共形不变性的拓扑事实：实轴与圆周共形等价（只差一个点），故任何作为模参数函数的奇点集在两者上相同。作者随即点出真正薄弱环节——(3.16) 的 $U$ 由大 $\sigma$ 极限定义，而红外极限是 $e\to\infty$；两个极限是否交换无法在有限体积的半经典框架内判定。这正是后文转向“准拓扑量”的动机：只研究那些与 $e$、与 Kähler 度规的 D 项形变无关的量，便可绕开 $e\to\infty$ 与半经典近似的冲突。]

无论如何，在半经典层面，Calabi--Yau 与 Landau--Ginzburg 就像液体与气体一样显得是不同的“相”；作者将沿用这一方便而非正式的用语。

\subsection{应用}

要点：本小节要确定在有限规范耦合 $e$ 下定义的线性 $\sigma$ 模型与共形不变的非线性 $\sigma$ 模型之间，哪些量是可靠的、可计算的；答案的核心是“准拓扑量”，以及由此引出的指标、椭圆亏格与扭转模型。

无论 $r$ 为正还是为负，这里所研究的线性 $\sigma$ 模型都是超可重整的（super-renormalizable），而非共形不变的。即使在经典层面，要把 $r\gg0$ 的理论约化为 Calabi--Yau 非线性 $\sigma$ 模型，或把
% ==== 原文 PDF 第 30 页（印刷页 30）====
\label{op:30}
% [跨页接续: 原文第 29 页/第 30 页，句子跨页]
$r\ll0$ 的理论约化为 Landau--Ginzburg 模型，我们都必须取规范耦合 $e$ 趋于无穷的极限。

但我们的论证依赖于 $e<\infty$。再者，$e\ll1$ 才真正是模型处于有效控制之下的区域，此时半经典论证是可靠的。

% [Tutor Note: 两个相的几何来源见 §3.1：解辅助场方程得 $D=-e^{2}\bigl(\sum_i Q_i|\phi_i|^{2}-r\bigr)$，故 $D=0$ 要求 $\sum_i|s_i|^{2}-n|p|^{2}=r$（式 (3.6)）。$r>0$ 时 $s_i$ 不能全为零，对 $U(1)$ 取商得 $\mathbf{CP}^{n-1}$，再令 $G=0$ 得超曲面 $X$，凯勒类正比于 $r$；$r<0$ 时须 $p\neq0$、$s_i=0$，残余的 $\mathbb{Z}_n$ 规范对称性给出 Landau--Ginzburg 轨形。两种描述是同一模型在不同 $r$ 区间的低能极限。]

如果我们确实在 $r\gg0$ 下取 $e\to\infty$，就得到靶空间为超曲面 $X$ 的 Calabi--Yau 非线性 $\sigma$ 模型。$X$ 上的凯勒度规 $g$（以及相应的凯勒形式 $\omega$）是由到 $\mathbf{CP}^{n-1}$（配以其标准度规）的嵌入诱导而来的。已知 $X$ 上的这个度规并不给出共形不变的模型。人们相信（对足够大的 $r$）在 $X$ 上存在唯一的凯勒度规 $g'$，其凯勒形式 $\omega'$ 与 $\omega$ 上同调等价，且给出共形不变的 $\sigma$ 模型。度规为 $g$ 与 $g'$ 的两个模型相差一个如下形式的超空间相互作用
\begin{equation}
\tag{3.24}
\int \mathrm{d}^{2}x\,\mathrm{d}^{4}\theta\; T,
\end{equation}
其中函数 $T$ 满足 $\omega-\omega'=-i\partial\bar\partial T$。

% [Tutor Note: $\omega-\omega'=-i\partial\bar\partial T$ 正是凯勒势的移动：若 $\omega=i\partial\bar\partial K$、$\omega'=i\partial\bar\partial K'$，则 $T=K-K'$（在差一个 $\omega$ 的调和部分的意义下）。因此 (3.24) 是把“错误”度规 $g$ 换成共形不变度规 $g'$ 的整套超对称完成，属于 $D$ 项；它不改变任何 $D$ 项不变量，但会改变度规本身。注意符号约定：本文 $\omega$ 与凯勒势的关系固定了 $T$ 的整体符号，换用相反的 $\omega$ 约定时 $T\to-T$。]

% [Tutor Note: 超势单项式集合与射影空间权重之间的对应在此已现雏形。按 §3.1，$W=P\cdot G(S_1,\ldots,S_n)$，$G$ 是 $n$ 次齐次多项式，规范荷为 $Q_{S_i}=1$、$Q_{P}=-n$；规范不变性要求 $G$ 的每个单项式都恰为 $n$ 次齐次，其指数组 $(\alpha_1,\ldots,\alpha_n)$ 满足 $\sum_i\alpha_i=n$，而 $G=0$ 定义了环境空间 $\mathbf{CP}^{n-1}$ 中的超曲面 $X$。$\sum_iQ_i=0$ 等价于 $\deg G=n$，这正是 $X$ 为 Calabi--Yau 的条件。横向性（transversality）条件 $\partial G/\partial S_1=\cdots=\partial G/\partial S_n=0$ 只有公共根 $S_i=0$，它等价于 $X$ 光滑；这是本节把 $X$ 上诱导度规（虽非共形不变）当作“错误但光滑”的凯勒度规、并在其上定义椭圆亏格与指标的前提。$G$ 的单项式集合一旦固定，$X$ 的复结构模就由这些单项式的系数给出；§5.1 把这个对应推广到加权射影空间中的超曲面，那里超势的单项式集合由环境空间的权 $w_i$ 与次数 $d$ 决定，横向性条件的形式完全平行。]

因此，即使取 $e=\infty$，这里所考虑的线性 $\sigma$ 模型也不被期望与共形不变的非线性 $\sigma$ 模型相重合。二者相差 (3.24) 这一项，而且函数 $T$ 是未知的。那么，我们怎样才能学到与共形不变模型相关的任何东西呢？

一种答案是：$T$ 耦合与有限 $e$ 的效应一样，被认为在重整化群的技术意义下是“无关的”（irrelevant）；也就是说，这些效应被期望在红外消失。这在理论上是重要的——例如它出现在关于 Landau--Ginzburg/Calabi--Yau 对应关系的启发式讨论中——但难以作为计算的基础。

一个更实用的答案依赖于如下事实：二维的 $N=2$ 理论拥有丰富的量，它们在拉格朗日量上添加形如 $\int \mathrm{d}^{4}\theta(\cdots)$ 的项（其中 $\cdots$ 是规范不变算符）时保持不变。这类量包括 $\mathrm{Tr}(-1)^{F}$、椭圆亏格，以及两个扭转拓扑场论（$A$ 模型与 $B$ 模型）和一个迄今研究甚少的杂合体——半扭转模型（half-twisted model）——的可观测量。我们把这些量称为准拓扑量（quasi-topological quantities）。任何这样的量在拉格朗日量上添加诸如 (3.24) 的项时都不变，因此可以用 $X$ 上那个“错误的”凯勒度规来计算它们。

% ==== 原文 PDF 第 31 页（印刷页 31）====
\label{op:31}
% [跨页接续: 原文第 30 页/第 31 页；上一段末句（"这类量包括……"）自第 30 页末跨至第 31 页首，按语义合并译出；本锚点起的以下文字均在第 31 页]
此外，一旦研究者安于主要研究准拓扑量，就容易看出：我们在 $e$ 依赖上的困难可以按本质上相同的方式解决。拉格朗日量中依赖于 $e$ 的项是
\begin{equation}
\tag{3.25}
-\frac{1}{4e^{2}}\int \mathrm{d}^{2}y\,\mathrm{d}^{4}\theta\; \bar\Sigma\Sigma,
\end{equation}
因此准拓扑量与 $e$ 无关。对 Fayet--Iliopoulos 项与 theta 项的依赖可以类似地分析，结果是：某些准拓扑项在这些耦合下不变，另一些则全纯地依赖于 $t$。

% [Tutor Note: 规范动能项在超空间中就是 $\int \mathrm{d}^{4}\theta\,\bar\Sigma\Sigma$ 型的 $D$ 项，其系数 $1/4e^{2}$ 不进入任何“超荷恰为 $\theta$ 全微分”的积分；扭转之后这类项成为 $\{Q,\cdots\}$，故其真空期望与关联函数与 $e$ 无关。这恰好化解了第 29 页末的次序问题：可以先在 $e\ll1$ 的半经典区精确计算这些量，再取 $e\to\infty$，结果不变。但要强调，这一论证只对准拓扑量成立，不适用于一般的 $F$ 项或依赖于度规的量。]

\subsubsection*{指标}

最简单的准拓扑量就是各种算符的指标（在圆周上的正则量子化中）。如前所述，设 $H$ 与 $P$ 为理论的哈密顿量与动量，并令 $H_{\pm}=(H\pm P)/2$。可以找到超荷 $Q_{\pm}$ 与 $\bar Q_{\pm}$ 的一个线性组合 $Q$，使得 $Q$ 是自伴的且 $Q^{2}=H$。令 $(-1)^{F}$ 为把费米子数目按模 2 计数的算符，于是它与 $Q$ 反交换。对象 $\mathrm{Tr}(-1)^{F}e^{-\beta H}$ 由标准论证与 $\beta$ 无关（该论证涉及在非零 $H$ 的态之间由乘以 $Q$ 给出的配对）。这个量是 $Q$ 的指标（更确切地说，是 $Q$ 中把 $(-1)^{F}$ 偶的态映到奇的态的那一部分的指标）。它通常就简写为
\begin{equation}
\tag{3.26}
\mathrm{Tr}(-1)^{F}.
\end{equation}

% [Tutor Note: 推导补足。$Q^{2}=H$ 且 $Q$ 自伴，故 $Q$ 与 $H$ 对易、保持每个能量本征空间，并在 $H>0$ 的本征空间上可逆（逆为 $Q/H$）；于是每个这样的本征空间内玻色态与费米态的维数相等（更具体地，$Q$ 的本征值成对出现为 $\pm\sqrt{\lambda}$），对迹的贡献逐对抵消。只有 $H=0$ 的子空间不被配对，故 $\mathrm{Tr}(-1)^{F}e^{-\beta H}=n_{B}^{0}-n_{F}^{0}$，与 $\beta$ 无关。前提是 $H$ 在临界值以下的谱离散，否则这一配对论证需要额外的谱假设。]

这个指标与 $r$ 和 $\theta$ 无关（只要远离 $r=\theta=0$ 处的奇点），因为在某个临界值以下 $H$ 的谱是离散的；因此它在 Landau--Ginzburg 与在 Calabi--Yau 中具有相同的值。

% ==== 原文 PDF 第 32 页（印刷页 32）====
\label{op:32}
% [跨页接续: 原文第 31 页/第 32 页；上一句（"这个指标与 $r$ 和 $\theta$ 无关……"）自第 31 页末跨入第 32 页，本锚点起的以下文字均在第 32 页]
利用理论中左行与右行的 $R$ 对称性，可以构造一个类似但精细得多的不变量。令 $W_{L}=\oint J_{L}$ 与 $W_{R}=\oint J_{R}$ 为 $R$ 荷；并令
\begin{equation}
\tag{3.27}
\begin{aligned}
(-1)^{F_{R}}&=\exp(i\pi W_{R})\\
(-1)^{F_{L}}&=\exp(i\pi W_{L}).
\end{aligned}
\end{equation}
于是 $(-1)^{F}=(-1)^{F_{L}}(-1)^{F_{R}}$。令 $Q_{R}=(Q_{+}+\bar Q_{+})/2$，则 $Q_{R}$ 与 $(-1)^{F_{R}}$ 反交换、与 $(-1)^{F_{L}}$ 交换，且 $Q_{R}^{2}=H_{+}$。由标准论证（在 $H_{+}\neq0$ 的态之间由乘以 $Q_{R}$ 给出的配对），量
\begin{equation}
\tag{3.28}
\mathrm{Tr}(-1)^{F_{R}}q^{H_{-}}e^{i\theta W_{L}}\exp(-\beta H_{+})
\end{equation}
与 $\beta$ 无关；我们把它记为
\begin{equation}
\tag{3.29}
F(q,\theta)=\mathrm{Tr}(-1)^{F_{R}}q^{H_{-}}e^{i\theta W_{L}}.
\end{equation}

% [Tutor Note: 易错预警。(3.28) 与 (3.29) 的迹只对右行费米子数取反（$F_{R}$），而空间积分给出的是 $W_{L}=\oint J_{L}$（左行 $R$ 荷）；若把 $F_{L}$ 与 $F_{R}$、或把 $W_{L}$ 与 $W_{R}$ 对调，得到的迹一般不同，甚至可能因配对而为零。还应区分两个记号：$W_{L},W_{R}$ 是 $R$ 荷算符（积分后的荷），而 (3.30) 中的 $J_{L},J_{R}$ 是相应的流；$\oint$ 表示沿圆周空间截面的积分。]

这个量可以解释为 $Q_{R}$ 的指标，它本质上就是该模型的椭圆亏格 [26--28]；它具有有趣的模性质，这些性质可以由其路径积分表示来建立。鉴于临界值以下 $H_{+}$ 的谱是离散的，$F(q,\theta)$ 在超对称参数变动下不变，因此它是拓扑不变量，而且在 Calabi--Yau 与 Landau--Ginzburg 中具有相同的值。（在 Landau--Ginzburg 情形，这个函数此前尚未被研究过；它具有相当有趣的性质，将在别处探讨。）

% [Tutor Note: 推导补足。与 (3.26) 的配对论证平行：$Q_{R}^{2}=H_{+}$，故在 $H_{+}>0$ 的子空间上 $Q_{R}$ 给出成对消去，于是取 $\beta\to\infty$ 或 $\beta\to0$ 结果相同，$F(q,\theta)$ 与 $\beta$ 无关。$q$ 通常取 $q=e^{2\pi i\tau}$，$\tau$ 为世界面环面的模参数；该函数在 $\mathrm{SL}(2,\mathbb{Z})$ 下的模性质来自环面上的路径积分，这也说明为什么需要 $(-1)^{F_{R}}$ 无反常（即靶空间为自旋流形），否则不同自旋结构上的配分函数不能拼成一个模形式。]

如果 $\sum_{i}Q_{i}$ 是偶数但不为零，则只要把 $\theta$ 限制为 $2\pi\bigl(\sum_{i}Q_{i}\bigr)^{-1}$ 的整数倍（对应于无反常的离散 $R$ 对称性），并适当计入大 $\sigma$ 处的真空态，$F(q,\theta)$ 仍然是一个不变量。需要 $\sum_{i}Q_{i}$ 为偶数是为了保证 $(-1)^{F_{R}}$ 无反常；在低能下，它导致 $r\gg0$ 的 $\sigma$ 模型的靶空间是一个自旋流形，而这是定义椭圆亏格时的一个自然要求。

% [Tutor Note: 对本节所关心的模型，$\sum_{i}Q_{i}=n\cdot1+(-n)=0$（$S_{i}$ 的荷为 $1$，$P$ 的荷为 $-n$，见 §3.1），零是偶数，故此处“偶数但不为零”的一般性限制对本模型自动满足，无需把 $\theta$ 离散化。$\sum_{i}Q_{i}=0$ 正是 §2 的反常抵消（无反常 $R$ 对称性）条件，也是 $B$ 模型无反常的条件（见后文第 33 页）；而低能靶空间的自旋性则对应 $X$ 的第一陈类为零。]

% ==== 原文 PDF 第 33 页（印刷页 33）====
\label{op:33}

\subsubsection*{延拓到欧氏世界面；扭转}

本文至此，二维世界面一直取洛伦兹号差。对于余下的应用，欧氏世界面更为自然。我们通过设 $y^{0}=-iy^{2}$ 解析延拓到欧氏号差的世界面；线元 $\mathrm{d}s^{2}=-(\mathrm{d}y^{0})^{2}+(\mathrm{d}y^{1})^{2}$ 变为 $\mathrm{d}s^{2}=(\mathrm{d}y^{1})^{2}+(\mathrm{d}y^{2})^{2}$。

这些线性 $\sigma$ 模型不仅不是拓扑不变的；它们甚至不是共形不变的，或许在红外极限下除外。然而，正如任何二维 $N=2$ 超对称模型一样，可以通过一个简单的扭转程序构造拓扑场论 [29,24,30]。这涉及把能量动量张量替换为
\begin{equation}
\tag{3.30}
T_{\alpha\beta}\to T_{\alpha\beta}+\frac{1}{4}\Bigl(\epsilon_{\alpha}{}^{\gamma}\partial_{\gamma}(J_{R\beta}\pm J_{L\beta})+\alpha\leftrightarrow\beta\Bigr).
\end{equation}

依符号的选择，得到两个模型。对 $+$ 号与 $-$ 号，分别得到所谓的 $B$ 模型与 $A$ 模型。$B$ 模型除非 $\sum_{i}Q_{i}=0$，否则有反常，因此我们对它的讨论限于该情形。$A$ 模型没有这种限制（但当 $\sum_{i}Q_{i}=0$ 时 $A$ 模型具有非常特殊的性质，因为只有那时瞬子求和才是无穷的）。

% [Tutor Note: 扭转的实质是把能量动量张量作一个“改进”：加上流的散度项（这里用二维反对称张量 $\epsilon_{\alpha}{}^{\gamma}\partial_{\gamma}$ 写成），使得某一对超荷的洛伦兹自旋由 $\pm\frac12$ 变为 $0$，从而在任意黎曼面上都有定义。$\pm$ 的选择决定保留哪一对标量超荷：按本文约定，$B$ 模型保留 $\bar Q_{+},\bar Q_{-}$，$A$ 模型保留 $Q_{-},\bar Q_{+}$（下一段）。易错预警：不同文献对 $T_{\alpha\beta}$ 的符号与 $\epsilon$ 的约定（$\epsilon_{12}=-1$ 或 $+1$）不同，扭转公式的形状随之改变；核对时不要套用别的文献的 $\pm$ 约定。]

扭转程序的要点在于：若洛伦兹变换是用修改后的能量动量张量来定义的，则某些超荷会具有零洛伦兹自旋。对 $B$ 模型，这些超荷是 $\bar Q_{+}$ 与 $\bar Q_{-}$；对 $A$ 模型，则是 $Q_{-}$ 与 $\bar Q_{+}$。由于这些量是洛伦兹不变的，当模型被建立在任意亏格 $g$ 的黎曼面 $\Sigma$ 上时，它们仍然是模型的对称性。因此，无论是 $A$ 模型还是 $B$ 模型，建立在任意曲面 $\Sigma$ 上时都具有一个 $(0|2)$ 维的对称超群 $F$。

令 $Q=\bar Q_{+}+\bar Q_{-}$（$B$ 模型）或 $Q=Q_{-}+\bar Q_{+}$（$A$ 模型）；则在两种情形都有 $Q^{2}=0$。一个自然的想法是把 $Q$ 当作某种 BRST 算符，把物理态看作 $Q$ 的上同调类，把可观测量看作 $Q$ 不变顶点算符的关联函数。可以证明（如同经典 Hodge 理论中那样），BRST 上同调类具有 $F$ 不变的代表元；这带来一些有用的简化。

% [背景补充: 经典 Hodge 理论中，de Rham 上同调类在每一类中都有唯一的调和代表元；这里所类比的结论是：在 $Q$ 的上同调类中总可以选取被 $F$ 保持的代表元，于是路径积分可以把 $F$ 当作对称性来用。]

% ==== 原文 PDF 第 34 页（印刷页 34）====
\label{op:34}
% [跨页接续: 原文第 33 页/第 34 页；上一句（"可以证明（如同经典 Hodge 理论中那样）……"）自第 33 页末跨入第 34 页，本锚点起的以下文字均在第 34 页]
在两个扭转模型中的任何一个，都有 $T_{\alpha\beta}=\{Q,\Lambda_{\alpha\beta}\}$（对某个 $\Lambda_{\alpha\beta}$）；这一条件保证 $Q$ 不变算符的关联函数与度规无关。于是扭转模型成为拓扑场论；我们称其可观测量为拓扑关联函数。正如 [24] 中详细解释的，$A$ 模型与 $B$ 模型的拓扑可观测量作为特例包含了未扭转模型的 $ac$ 与 $cc$ 手征环，而这些手征环决定了未扭转模型所定义的弦紧化的低能 Yukawa 耦合。

拉格朗日量中任何形如 $\int \mathrm{d}^{2}y\,\mathrm{d}^{4}\theta\cdots$ 的微扰都可以写成 $\{Q,\cdots\}$，因此拓扑可观测量在此类微扰下不变。特别地，拓扑可观测量与规范耦合 $e$ 无关。对其他耦合，基本规则如下。在 $B$ 模型中，拉格朗日量中任何可以写成 $\int \mathrm{d}\bar\theta^{+}\cdots$ 或 $\int \mathrm{d}\bar\theta^{-}\cdots$ 的项都是 $\{Q,\cdots\}$；而在 $A$ 模型中，任何可以写成 $\int \mathrm{d}\theta^{-}\cdots$ 或 $\int \mathrm{d}\bar\theta^{+}\cdots$ 的项都具有这一性质。

考察超势耦合的定义
\begin{equation}
\tag{3.31}
L_{W}=\int \mathrm{d}^{2}y\,\mathrm{d}\theta^{+}\mathrm{d}\theta^{-}\;W(\Phi_{i})\big|_{\bar\theta^{+}=\bar\theta^{-}=0}+\int \mathrm{d}^{2}y\,\mathrm{d}\bar\theta^{-}\mathrm{d}\bar\theta^{+}\;\overline{W(\bar\Phi)}\big|_{\theta^{+}=\theta^{-}=0},
\end{equation}
即可看出：$A$ 模型的可观测量与超势无关，而 $B$ 模型的可观测量随 $W$ 全纯地变化。

% [Tutor Note: 推导补足。上一段给出的规则之所以成立，是因为扭转后 $Q$ 中含 $\bar Q_{\pm}$（$B$ 模型）或 $Q_{-},\bar Q_{+}$（$A$ 模型），它们在超场上的作用等价于（在相应手征坐标中）对某一个 $\theta$ 的平移/求导；于是凡在相应 $\theta$ 上为全微分、只做单个 Grassmann 积分的项，都可以写成 $\{Q,\cdots\}$。把这套规则套到 (3.31)：第一项只含手征 $\theta^{+},\theta^{-}$ 的积分，第二项只含反手征 $\bar\theta^{-},\bar\theta^{+}$；在 $B$ 模型中第二项（反手征）成为 $\{Q,\cdots\}$，故只有 $W$ 的手征部分（全纯依赖）留下；在 $A$ 模型中第一项（手征）成为 $\{Q,\cdots\}$，故 $A$ 模型与 $W$ 无关。原文的 "we see that" 隐含的正是这一步。]

类似地，回顾 Fayet--Iliopoulos 与 theta 耦合的定义
\begin{equation}
\tag{3.32}
L_{D,\theta}=\frac{it}{2\sqrt{2}}\int \mathrm{d}^{2}y\,\mathrm{d}\theta^{+}\mathrm{d}\bar\theta^{-}\;\Sigma\big|_{\theta^{-}=\bar\theta^{+}=0}-\frac{i\bar t}{2\sqrt{2}}\int \mathrm{d}^{2}y\,\mathrm{d}\theta^{-}\mathrm{d}\bar\theta^{+}\;\bar\Sigma\big|_{\theta^{+}=\bar\theta^{-}=0},
\end{equation}
即可看出：$B$ 模型的拓扑可观测量与 $r$、$\theta$ 无关，而 $A$ 模型的拓扑可观测量全纯地依赖于 $t$。

% [Tutor Note: 物理图像。(3.32) 的两个系数互为复共轭，$t=ir+\theta/2\pi$（式 (2.28)）是复化的 FI 参数。按上一段的一般规则：$B$ 模型中形如 $\int\mathrm{d}\bar\theta^{\pm}\cdots$ 的项是 $\{Q,\cdots\}$，故 (3.32) 的两项（分别含 $\mathrm{d}\bar\theta^{-}$ 与 $\mathrm{d}\bar\theta^{+}$）都成为 $\{Q,\cdots\}$，$B$ 模型与 $t$ 的实部 $r$ 和虚部 $\theta$ 都无关；$A$ 模型中形如 $\int\mathrm{d}\theta^{-}\cdots$ 或 $\int\mathrm{d}\bar\theta^{+}\cdots$ 的项是 $\{Q,\cdots\}$，故只有第二项（含 $\mathrm{d}\theta^{-}\mathrm{d}\bar\theta^{+}$）成为 $\{Q,\cdots\}$，留下系数为 $it$ 的第一项，$A$ 模型因而只全纯地依赖 $t$。这与镜像对称的标准图像一致：$B$ 模型依赖复结构（超势 $W$），$A$ 模型依赖复化凯勒类 $t$。由此，$t$ 的解析延拓就是把 Calabi--Yau 延拓到 Landau--Ginzburg 的具体实现。]

% [跨页接续: 本片止于第34页末完整句，续见 s03d]
